\subsubsection{SEM Frequently Asked Questions}\label{SEC:SEM:FAQ}

\begin{QandA}
%
  \item[Limitations on the physical inflow boundary: must it be planar?]
    Yes, the inflow boundary must be planar. The algorithm that interpolates the
    SEM profile onto the inflow boundary points uses wall distance and the wall
    normal to rotate and interpolate the profile. The wall distance part would
    cause problems with a non-planar face.
%
  \item[Limitations on the physical inflow boundary: must it be constant-$x$?]
    Yes. In its current form, the SEM can only be used on a constant-$x$ inflow
    boundary. One reason for this is that the synthetic eddies are only
    propagated in $x$. 
%
  \item[Can the SEM be used to add freestream turbulence to a \quoterm{farfield}
    subsonic inflow boundary, \thatis one that is not adjacent to a wall
    boundary?]
    No, not at the moment. The SEM initialization algorithm is dependent on
    finding the nearest wall boundary point for each each inflow boundary point.
    This is used to compute wall distance and get the wall normal at the wall
    point, which are both used to interpolate and rotate the SEM profile data to
    get the imposed flow profile at the inflow point. Attempting to use the SEM
    in this way would almost certainly cause either a crash via segmentation
    fault or (if it somehow got past the initialization) a badly imposed
    flow profile at the inflow. We are hoping to add this capability to the SEM
    in the near future.
%
  \item[Is the virtual volume extruded normal to the physical inflow boundary or
    along some velocity vector (based on inflow profile)?]
    Neither. It takes the min\slash max of each coordinate direction over all
    the boundary points on the inflow plane and respectively subtracts\slash
    adds the SEM length scale to each value to get the extents of the virtual
    box.
%
  \item[Are the eddies convected normal to the physical inflow boundary or along
    some velocity vector (based on inflow profile)?]
    The eddies are currently convected in only the $x$ direction.
%
  \item[Is the sensory distance between an eddy and face point done in 3D (as a
    sensory spherical volume) or 2D (within the physical inflow plane)?]
    A tent function is used to get a contribution in each direction,
    then all three contributions are multiplied together to get a total
    fluctuating value. There is an input to treat the sensory distance of each
    eddy as a sphere,
   %\\% Prevent next word from entering right margin.
    \inp[nmlSEM::UseSphericalEddyVolume]{Use\_Spherical\_Eddy\_Volume},
    but this only affects how the eddy volume factors in to computing the total
    number of eddies within an individual SEM box and has no effect when
    computing the contribution of each eddy to the velocity fluctuations for the
    boundary condition.
%
  \item[What characteristic information feeds upstream?]
    The SEM does not use any information from the interior to set the boundary
    states.
    %
    However, as discussed in \sref{SEC:SEM:usage:subsonic}, any points on an
    SEM inflow boundary that are outside of the provided boundary layer profile
    are automatically separated into a standard subsonic inflow boundary
    that does permit outgoing characteristic information.
%
  \item[For subsonic flow, does SEM overprescribe the flow conditions?]
    Yes, the SEM ends up prescribing all flow variables for the boundary state
    and does not use the interior state at all. This ultimately results in a
    \quoterm{supersonic} or \quoterm{fixed} inflow type boundary condition.
    However, this seems to cause little to no numerical issues, so we just tend
    to ignore this and take the win.
%
  \item[Should the entire inflow boundary be prescribed SEM or only some
    portion?]
    It does not really matter. The parts of the inflow boundary outside the SEM
    profile height will be separated out to a new BC. If you prescribe the
    entire inflow boundary as SEM, it will result in two BC's: one for the SEM
    and one for a subsonic inflow BC for those faces outside the SEM profile
    height. Manually separating out the SEM and non-SEM parts of the inflow
    boundary will result in three BC's: the same two as before plus one for the
    manually separated out boundary faces. The difference in computational cost
    between these two should be negligible. 
%
  \item[Is more than one SEM permitted?]
    Yes, multiple SEM boundaries can be specified. It is possible, but unlikely,
    that the virtual boxes of two or more SEM BC's overlap slightly, but each
    SEM and its corresponding set of boundary faces are completely
    self-contained, meaning there is no influence from a different SEM.
%
  \item[How are corner flows handled?]
    The distance variable in the SEM profile is height off the wall. The nearest
    wall point for each face point is found, and the SEM profile is interpolated
    to that face point using the normal at the wall point with the help of
    rotation transformation matrices to orient the flow velocities and stresses
    relative to the wall. For an inflow boundary containing a corner where two
    walls meet, this interpolation results in a sort of pinch-point of the flow
    profiles at the corner. This may not be the best 2D profile to specify at
    the inflow, but it seems to result in a reasonable corner flow by the time
    the unsteadiness transitions into a physical turbulent boundary layer.
%
  \item[How much streamwise distance is required to achieve realistic
    turbulence?]
    If using an SEM convective velocity and length scale within the recommended
    values provided in \sref{SEC:SEM:usage:recommendations}, it could take
    anywhere from 10-15 \lBC to develop realistic turbulence statistics in line
    with the SEM profile, with 12 \lBC being a good starting estimate.
%
  \item[If the inflow profile represents the desired goal at some reference
    station, but the inflow boundary must be placed upstream for realistic
    turbulence to develop, should the profile height be scaled to account for
    boundary layer growth by the time it reaches the reference station?]
    The height of the boundary layer once it has achieved realistic turbulence
    should be reasonably close to the SEM profile height. Therefore, based on
    the previous question, the inflow boundary should be placed 12-15 \lBC
    upstream of the location you would like the SEM profile exist in the
    simulation flow field.
%
  \item[Is it best to only specify the boundary layer profile up to but not
    beyond the boundary layer edge?]
    Yes. By providing a profile file with data that only extends to the
    edge of the boundary layer, more boundary points will be
    separated out to the new subsonic inflow boundary condition. The traditional
    subsonic inflow BC is more efficient than the SEM BC. However, the SEM BC is
    quite efficient, so the difference between a full subsonic inflow BC and a
    full SEM inflow BC is probably less than 10\% computational cost.
%
  \item[How to specify which inflow(s) use SEM and which do not?]
    An SEM boundary is specified by setting
    \BCinpBC[nmlBC::BCInput/SetBCValue]{Set\_BC\_Value}=5
    in the \inp[nmlBC]{nml\_bc} namelist input block.
%
  \item[What input specifies the file name for the SEM profile?]
    \BCinpBC[nmlBC::BCInput/SEM/ProfileFile]{SEM\%Profile\_File}%
    =\texttt{"}[Path to profile file]\texttt{"}
    in the \inp[nmlBC]{nml\_bc} namelist input block.
%
 %\item[How is the virtual volume constructed from the physical inflow
 %  boundary?]
 %  See \sref{SEC:SEM:theory:VirtualBox}
%
 %\item[How are turbulent fluctuations to the thermo variables ($\rho,P,T$)
 %  handled?]
 %  See \sref{SEC:SEM:theory:ComputeBC}.
%
 %\item[How does SEM handle flow beyond the extent of the provided profile?]
 %  See \sref{SEC:SEM:usage:subsonic} and \sref{SEC:SEM:usage:example}.
%
  \item[Is the SEM profile 1D or 2D?]
    The SEM profile is specified as a 1D profile as described in
    \sref{SEC:SEM:usage:profile}.
%
 %\item[What is the format of the SEM profile file?]
 %  See \sref{SEC:SEM:usage:profile:format}.
%
\end{QandA}

%###############################################################################

